\documentclass[11pt]{article}

\usepackage[margin=1in]{geometry}
\usepackage{amsmath,amssymb,amsthm,mathtools}
\usepackage{booktabs}
\usepackage{microtype}
\usepackage{setspace}
\usepackage{enumitem}
\usepackage[round,authoryear]{natbib}
\usepackage[colorlinks=true,allcolors=blue]{hyperref}

\newtheorem{proposition}{Proposition}
\newtheorem{definition}{Definition}
\newtheorem{remark}{Remark}

\newcommand{\E}{\mathbb{E}}
\newcommand{\R}{\mathbb{R}}
\newcommand{\dd}{\,\mathrm{d}}

\title{Automating Innovation:\\
Artificial Intelligence, Market Structure, and the Dynamics of Research}
\author{Oliver Pardo}
\date{August 2026}

\begin{document}

\maketitle

\begin{abstract}
This paper studies a dynamic economy in which artificial intelligence raises current
productivity and can also be used to produce better artificial intelligence.
Population grows exogenously, competitive final-good firms use quality-adjusted AI
services, and one integrated frontier developer supplies those services as a
monopolist and invests in AI research. Frontier improvements combine human and
automated-research services through a CES technology with an elasticity above one.
The automated share rises when the research wage increases relative to its unit
resource cost, and the human contribution can become asymptotically negligible.
During a human-dominated phase, capability growth remains semi-endogenous and tied
to population growth. When automated research dominates, capability growth is tied
instead to the growth of automated-research services. In the Cobb--Douglas production
benchmark, a constant-growth path exists only when diminishing returns in ideas
production dominate the feedback between output and automated research. With a
general CES between production labor and AI services, complementarity, constant
shares, and strong substitution produce different asymptotic growth regimes. The
same production elasticity separates the response of the real wage from that of the
production-labor share. Monopoly restricts current AI deployment but creates the
rents that give frontier capability a private value.
\end{abstract}

\onehalfspacing

\section{Introduction}
\label{sec:introduction}

Research has historically been constrained by the number, ability, and organization
of human researchers. Artificial intelligence may relax this constraint. The
relevant change would not occur when AI becomes useful in ordinary production, or
even when it automates a large share of existing jobs. It would occur when AI systems
can perform the research tasks required to produce better AI systems. At that point,
AI capability would become an input into its own law of motion, and innovation would
no longer be tied mechanically to the supply of human researchers.

The paper does not model the earlier transition from an economy without commercially
useful AI to one that uses AI in ordinary production. One can think of that
pre-AI period as an economy with conventional production and human research, in
which scientific knowledge, algorithms, and computing capacity accumulate before AI
services become economically useful. A formal model of that transition would require
an adoption margin and an explanation of how research is financed before it produces
current AI revenues. We instead begin after productive AI services already exist but
while frontier research remains primarily dependent on human researchers. The object
of the paper is the subsequent automation of AI research.

This possibility raises an economic question before it raises a forecasting question.
Suppose frontier improvements combine human researchers and automated-research
services. Firms then choose both the scale and the composition of research. When the
two inputs are gross substitutes, the automated share increases if the human research
wage rises relative to the unit resource cost of automated research. Automation is
therefore a continuous equilibrium transition rather than a date imposed on the
model. Human researchers can become asymptotically inessential, but substitution
above one does not guarantee that outcome: the wage must also rise relative to the
cost of automated research.

Population growth is important during the human-dominated phase. In standard R\&D-based growth
models, new ideas depend on a growing but scarce human research input
\citep{romer1990,jones1995,aghionhowitt1992}. With diminishing returns to the stock of
ideas, a constant research-employment share implies semi-endogenous growth: the
long-run rate of innovation is pinned down by the growth of the research labor force.
Our benchmark preserves this logic while human research dominates. If population grows at rate
$n$, the capability growth rate is $\eta n/(1-\phi)$; without population growth,
positive asymptotic growth disappears when $\phi<1$.

Research automation can alter this result. Effective research labor can expand
through additional automated-research services rather than through population growth.
When those services dominate, capability growth is tied to their equilibrium growth
rate. Innovation can accelerate if higher output finances enough additional
automated research, but diminishing returns to ideas and research effort can preserve
a constant-growth path. The outcome depends on the elasticities in production and
research and on the allocation of output to research.

Recent discussions often move too quickly from this mechanism to one of two
conclusions. One view treats an intelligence explosion as the natural implication of
automating AI research. The other assumes that physical bottlenecks and diminishing
returns will necessarily preserve familiar growth rates. Neither conclusion follows
without specifying the ideas production function. Scenario exercises such as
\citet{kokotajlo2025} make the feedback mechanism concrete, but rapid takeoff depends
on several jointly necessary assumptions about research automation, algorithmic
progress, compute, and organizational capacity. An economic model should expose those
assumptions separately.

Market structure determines both current use and the private return to innovation.
Frontier AI requires large fixed investments and may exhibit economies of scale and
scope \citep{korinekvipra2025,gans2024}. A monopolist restricts the current supply of
AI services but captures operating rents that raise the private value of better
capability. We therefore begin with one integrated developer that both supplies
quality-adjusted AI services and invests in the frontier. Competitive final-good
firms combine capital with a CES composite of production labor and AI services. The
elasticity of substitution inside this composite governs adoption and the demand
elasticity faced by the monopolist.

The same elasticity separates the response of the real wage from the response of the
production-labor share. An expansion of AI services reduces this share when AI and
production labor are gross substitutes, but the wage falls only when substitution is
sufficiently strong. There is therefore an intermediate region in which production
workers receive a smaller share of output even though their wage rises. This
distinction is important for the distributional analysis: a declining labor share is
not, by itself, evidence that the real wage has fallen. In the full model, the
aggregate labor share also depends on how workers are reallocated between production
and research.

The resulting welfare problem has two distinct margins. Static market power reduces
the current use of AI services. At the same time, consumer surplus and knowledge
spillovers imply that the developer may capture only part of the social value of
better AI and therefore invest too little in research. Reducing the markup can
correct the deployment distortion while also eroding the rents that finance
innovation. Price regulation and research policy must consequently be studied
together. The present version isolates this market-structure mechanism.

The toy model delivers five results. First, the automated share of research follows
a closed-form law driven by the research-substitution elasticity and the wage relative
to the resource cost of automated research. Second, population growth sustains
semi-endogenous capability growth while human research dominates. Third, when
automated research dominates, a constant-growth path exists only if diminishing
returns offset the feedback between output and research. Fourth, the production CES
separates a labor-bottlenecked regime, a constant-share benchmark, and a regime in
which the return to capital rises with capability. Fifth, the production structure
yields distinct substitution thresholds for the real wage and the
production-labor share.

The extended model retains the integrated monopolist and adds distributional
heterogeneity and richer capital and compute accumulation. It will compare a
laissez-faire monopoly with regulated access prices, research subsidies or prizes,
and policies that combine both instruments. The objective is not to select a policy
before solving the model. It is to identify when lower AI-service prices must be
paired with explicit support for frontier research.

The rest of the paper is organized as follows. Section~\ref{sec:literature} relates
the framework to endogenous growth, transformative AI, automation and labor income,
measurement of AI services, and industrial organization. Section~\ref{sec:toy}
develops the toy model.
Section~\ref{sec:extended} sets out the extended general-equilibrium model.
Section~\ref{sec:regulation} defines the regulatory scenarios.
Section~\ref{sec:conclusion} concludes.

\section{Related literature}
\label{sec:literature}

The paper connects five literatures that are usually studied separately.

\paragraph{Endogenous innovation and economic growth.}
The first literature makes ideas production an economic activity. In
\citet{romer1990}, nonrival knowledge and imperfect competition sustain investment in
new varieties. \citet{aghionhowitt1992} model growth through quality-improving creative
destruction. \citet{jones1995} shows that diminishing returns in the production of ideas
remove scale effects and tie long-run growth to the growth of the research labor force.
Our starting point is closest to the semi-endogenous formulation: human researchers
are initially the scarce input that limits innovation. The departure is that research
labor can itself be produced with AI and compute.

\paragraph{AI, automation, and the production of ideas.}
\citet{aghionjonesjones2019} examine AI as an automation technology and emphasize that
growth can remain constrained by tasks that are essential but difficult to automate.
\citet{acemoglu2025} also models AI at the task level and relates its aggregate
productivity effect to the share of tasks affected and the cost savings within those
tasks.
\citet{nordhaus2021} studies the substitutability conditions required for an economic
singularity. \citet{trammellkorinek2023} distinguish the automation of ordinary
production from the automation of R\&D and show that physical and technological
bottlenecks remain central even under transformative AI. \citet{jones2025rd} provides
a general framework in which AI's effect on research depends on the share of tasks it
can perform, its productivity in those tasks, and the strength of complementary
bottlenecks. Our model endogenizes the gradual composition of human and automated
research and embeds the ideas production function in a decentralized market
equilibrium.

\paragraph{Automation, wages, and the labor share.}
The effect of automation on labor income depends on both substitution and
productivity. \citet{acemoglurestrepo2019} distinguish the displacement effect of
automating existing tasks from the productivity and reinstatement effects that can
raise labor demand. In CES models, the unit-elasticity threshold determines whether
an increase in a factor reduces the income share of the other factor. This mechanism
underlies the account in \citet{karabarbounisneiman2014} of how cheaper investment
goods can reduce the labor share. A lower labor share, however, does not imply a lower
real wage. Using a nested CES and Korean firm and establishment data,
\citet{aumshin2024} estimate an elasticity of 0.6 between equipment and labor and an
elasticity of 1.6 between software and labor. The second estimate is a useful
empirical reference for our production elasticity $\sigma_{XL}$ because $X$ is a flow of
AI services rather than a stock of conventional equipment. It is not a direct
estimate of substitution between AI services and labor in the world economy.
\citet{autorkausik2026} establish more generally that automation can raise
wages while reducing labor's share in a competitive constant-returns economy.
\citet{yango2026} obtains, in a nested Cobb--Douglas--CES technology, the same two
thresholds that appear in our production block: the labor share falls when
$\sigma_{XL}>1$, while the wage rises when $\sigma_{XL}<1/\alpha$. Thus, for
$1<\sigma_{XL}<1/\alpha$, an expansion of AI simultaneously raises the wage and lowers
labor's share of income. \citet{trammellkorinek2023} place this ambiguity in the
broader transformative-AI literature: production automation can sharply reduce the
labor share, while the wage response continues to depend on technology and
accumulation.

We use the two-threshold result as a benchmark rather than claim it as a new
comparative static. Our departure is to make AI a flow of quality-adjusted services
supplied by an integrated monopolist and to make the underlying capability frontier
endogenous. Market power changes the equilibrium path of AI deployment, capital
accumulation, wages, and operating rents; frontier research determines how long the
economy remains in each production regime. This structure lets us study how service
regulation and research policy affect the wage and the labor share through different
static and dynamic channels.

\paragraph{AI services, inference compute, and quality adjustment.}
Using an existing model and producing a new one require separate measures of
inference and research compute. \citet{korinekmckelvey2026} combine inference compute
with intangible model capital and adjust the resulting output for quality.
\citet{demireretal2025} document prices per token within intelligence tiers, which
likewise implies that raw tokens are not homogeneous AI services. Our specification
$X_t=A_tU_t$ is an efficiency-units representation, not a technological identity:
$U_t$ measures the scale at which the current model is run, while $A_t$ converts
inference compute into quality-adjusted services. Task-based formulations such as
\citet{aghionjonesjones2019} and \citet{acemoglu2025} provide a richer alternative
when heterogeneity across uses of AI is central.

\paragraph{Market structure in AI.}
The industrial-organization literature emphasizes high fixed costs, scarce compute,
data, economies of scope, vertical integration, and feedback between model quality and
the surrounding ecosystem. \citet{korinekvipra2025} study how these forces can produce
concentration and tipping in foundation models. \citet{gans2024} emphasizes the role of
markets for training and input data in determining market power. Using model-usage
data, \citet{demireretal2025} document entry, falling prices, differentiation, and
turnover among leading models. The evidence is consistent with strong competition
today but does not rule out future concentration. Our extended model focuses on a
specific dynamic mechanism: a temporary lead becomes persistent when it lowers the
leader's cost of producing the next generation of frontier models.

Our baseline takes this concentration to its limiting case: one integrated frontier
developer supplies current services and conducts frontier research. Regulated access
prices are introduced as policy counterfactuals rather than as a second baseline
market structure. This formulation isolates the dynamic link between current
operating rents and the private incentive to improve the frontier.

\section{A toy model}
\label{sec:toy}

The toy model isolates the interaction between population growth, market power in
current AI services, and the automation of AI research. Time is continuous. There is
one final good, a representative dynasty, competitive final-good firms, and one
integrated frontier developer. The developer monopolistically supplies AI services
and invests in the next generation of AI. The initial economy already uses AI
services in production; the model does not describe their first commercial adoption.

\subsection{Household}

Population grows at the exogenous rate $n$:
\begin{equation}
    \dot N_t=nN_t,
    \qquad N_0>0,\quad n\geq0.
    \label{eq:population}
\end{equation}
The representative dynasty has preferences
\begin{equation}
    \int_0^\infty e^{-\rho t}N_t\log c_t\dd t,
    \qquad
    c_t\equiv\frac{C_t}{N_t},
    \qquad \rho>n,
    \label{eq:utility}
\end{equation}
where $C_t$ is aggregate consumption. Each household member supplies one unit of
labor inelastically. Labor is allocated between final production, $L_{Y,t}$, and
human AI research, $H_t$:
\begin{equation}
    L_{Y,t}+H_t=N_t.
    \label{eq:labor-clearing}
\end{equation}
The dynasty owns the capital stock and the frontier developer. Capital evolves as
\begin{equation}
    \dot K_t=I_t-\delta_KK_t.
    \label{eq:capital-law}
\end{equation}
Let $k_t\equiv K_t/N_t$. The per-capita budget constraint is
\begin{equation}
    \dot k_t
    =
    (r_t-\delta_K-n)k_t+w_t
    +\frac{\pi^X_t-w_tH_t-\xi M_t}{N_t}
    -c_t,
    \label{eq:household-budget}
\end{equation}
where $\pi^X_t$ is the developer's operating profit before research expenditure.
The household's Euler equation is
\begin{equation}
    \frac{\dot c_t}{c_t}=r_t-\delta_K-\rho.
    \label{eq:euler}
\end{equation}

\subsection{Final-good firms}

Competitive final-good firms combine physical capital with a composite of production
labor and AI services:
\begin{align}
    Y_t
    &=
    K_t^\alpha Z_t^{1-\alpha},
    \qquad \alpha\in(0,1),
    \label{eq:final-production}\\
    Z_t
    &=
    \left[
        (1-\omega)L_{Y,t}^{(\sigma-1)/\sigma}
        +\omega X_t^{(\sigma-1)/\sigma}
    \right]^{\sigma/(\sigma-1)},
    \qquad
    \omega\in(0,1),\quad \sigma>0.
    \label{eq:service-composite}
\end{align}
The elasticity of substitution between capital and the service composite is one.
The elasticity between production labor and AI services is $\sigma$. The
Cobb--Douglas composite is obtained by continuity when $\sigma=1$.

Define the contribution of AI services to the CES composite as
\begin{equation}
    s^X_t
    \equiv
    \frac{\omega X_t^{(\sigma-1)/\sigma}}
    {(1-\omega)L_{Y,t}^{(\sigma-1)/\sigma}
        +\omega X_t^{(\sigma-1)/\sigma}}.
    \label{eq:ai-composite-share}
\end{equation}
Let $p^X_t$ denote the price of AI services. Profit maximization gives
\begin{align}
    r_t
    &=\alpha\frac{Y_t}{K_t},
    \label{eq:capital-demand}\\
    w_t
    &=(1-\alpha)(1-s^X_t)\frac{Y_t}{L_{Y,t}},
    \label{eq:labor-demand}\\
    p^X_t
    &=(1-\alpha)s^X_t\frac{Y_t}{X_t},
    \label{eq:ai-demand}
\end{align}
and relative demand satisfies
\begin{equation}
    \frac{X_t}{L_{Y,t}}
    =
    \left[
        \frac{\omega}{1-\omega}
        \frac{w_t}{p^X_t}
    \right]^\sigma.
    \label{eq:relative-ai-demand}
\end{equation}

\subsection{Integrated AI monopolist}

The variable $U_t$ is a flow of \emph{inference compute}: standardized
computational resources used during period $t$ to run an existing model. It can be
measured in FLOPs, GPU-hours, or equivalent units of a reference processor per unit
of time; it is not a stock of chips. Let $\xi>0$ denote the common unit resource cost
of standardized compute. It summarizes the capital, energy, and other resources
needed to supply one unit of compute, independently of whether that unit is used for
inference or research.

The variable $X_t$ is a flow of quality-adjusted AI services used by final-good
firms. Its empirical counterpart could be benchmark-equivalent tokens, predictions
adjusted for accuracy, or equivalent task hours. This distinction follows the
measurement approach in \citet{korinekmckelvey2026}, who combine inference compute
with model capital and adjust inference output for quality, and is consistent with
the intelligence tiers documented by \citet{demireretal2025}.

The stock $A_t$ measures frontier AI capability. Under the normalization used here,
it also measures quality-adjusted services produced per unit of inference compute:
\begin{equation}
    X_t=A_tU_t.
    \label{eq:ai-services}
\end{equation}
Equation~\eqref{eq:ai-services} is an efficiency-units specification, not a
technological identity. Conditional on $A_t$, it imposes constant returns to
inference compute and compresses task coverage, quality, speed, and inference
efficiency into one scalar.

\begin{remark}[Capability and inference efficiency]
The same state variable $A_t$ governs broad frontier capability and the productivity
of inference compute. This restriction keeps the service-market and research
mechanisms connected and makes the benchmark tractable, but it need not hold
quantitatively. A natural extension is $X_t=B(A_t)U_t$, with a separate mapping from
capability to inference efficiency. We retain $B(A)=A$ in the baseline and keep this
distinction as a robustness extension.
\end{remark}

The integrated developer recognizes the inverse demand generated by the final-good
sector and solves
\begin{equation}
    \max_{X_t\geq0}
    \left\{
        p^X_t(X_t;K_t,L_{Y,t})X_t
        -\xi\frac{X_t}{A_t}
    \right\}.
    \label{eq:monopoly-service-problem}
\end{equation}
Its operating profit is
\begin{equation}
    \pi^X_t=p^X_tX_t-\xi\frac{X_t}{A_t}.
    \label{eq:operating-profit}
\end{equation}

\subsection{Frontier research}

New AI capability is produced with effective research services, $E_t$:
\begin{align}
    \dot A_t
    &=
    \chi A_t^\phi E_t^\eta,
    \qquad \chi>0,\quad \phi<1,\quad \eta\in(0,1),
    \label{eq:ideas}\\
    E_t
    &=
    \left[
        (1-\nu)H_t^{(\theta-1)/\theta}
        +\nu M_t^{(\theta-1)/\theta}
    \right]^{\theta/(\theta-1)},
    \qquad
    \nu\in(0,1),\quad \theta>1.
    \label{eq:effective-research}
\end{align}
Here, $M_t$ is a flow of standardized automated-research services. It includes
training, post-training, experiments, synthetic-data generation, and inference
performed by AI agents engaged in R\&D. The baseline normalizes one unit of these
services to one unit in the research aggregator and assigns it the common resource
cost $\xi$. A capability-dependent conversion from compute into automated-research
services is left for an extension.

The elasticity $\theta$ governs substitution between human and automated research.
The restriction $\theta>1$ makes the two inputs gross substitutes and allows
positive research output when $H_t=0$. It does not impose that human research
disappears. That outcome depends on the equilibrium path of the wage relative to
$\xi$. The exponent $\phi$ captures how the existing stock of knowledge affects the
productivity of both research inputs.

\begin{remark}[Capability-dependent research productivity]
An extension replaces $M_t$ inside \eqref{eq:effective-research} with
$\lambda(A_t)M_t$. Its effective unit cost is then $\xi/\lambda(A_t)$, and the
automation condition depends on $w_t\lambda(A_t)/\xi$ rather than $w_t/\xi$. In the
AI-dominated limit, a power function $\lambda(A)=\lambda_0A^\psi$ changes the
capability exponent from $\phi$ to $\phi+\psi\eta$. The baseline sets
$\lambda(A)=1$ to isolate substitution between human and automated research from
capability-dependent improvements in research efficiency.
\end{remark}

Let $q_t$ be the developer's private shadow value of capability and let
$F(A_t,H_t,M_t)$ denote the right-hand side of \eqref{eq:ideas}. Its current-value
Hamiltonian is
\begin{equation}
    \mathcal H^P_t
    =
    \pi^X_t-w_tH_t-\xi M_t
    +q_tF(A_t,H_t,M_t).
    \label{eq:private-hamiltonian}
\end{equation}

\subsection{Social planner}

Let $Q_t$ denote the social shadow value of capability. Unlike $q_t$, $Q_t$ includes
consumer surplus from a larger supply of AI services and knowledge benefits that the
developer cannot appropriate. The comparison between $Q_t$ and $q_t$ isolates the
knowledge and appropriability wedge.

\subsection{Equilibrium}

A decentralized equilibrium is a sequence
\[
    \{C_t,c_t,I_t,K_t,N_t,L_{Y,t},H_t,X_t,U_t,M_t,A_t\}_{t\geq0}
\]
and prices
\[
    \{r_t,w_t,p^X_t,q_t\}_{t\geq0}
\]
such that:
\begin{enumerate}[label=(\roman*)]
    \item the representative dynasty maximizes \eqref{eq:utility}, taking prices and
    developer profits as given;
    \item final-good firms maximize profits under
    \eqref{eq:final-production}--\eqref{eq:service-composite};
    \item the integrated developer chooses current AI services and frontier research
    to maximize its market value;
    \item population, capital, and capability follow
    \eqref{eq:population}, \eqref{eq:capital-law}, and \eqref{eq:ideas};
    \item labor satisfies \eqref{eq:labor-clearing}; and
    \item the aggregate resource constraint is
    \begin{equation}
        C_t+I_t+\xi(U_t+M_t)=Y_t.
        \label{eq:resource-constraint}
    \end{equation}
\end{enumerate}
There is no common physical-capacity constraint between $U_t$ and $M_t$. The two
uses of compute compete through the final-good resource constraint, not through an
additional exogenous cap.

\subsection{Solution}

\subsubsection{Wages and the labor share under an AI expansion}

Let
\begin{equation}
    \ell^Y_t\equiv\frac{w_tL_{Y,t}}{Y_t}
    =(1-\alpha)(1-s^X_t)
    \label{eq:production-labor-share}
\end{equation}
denote production labor's share of final output. The wage and this share need not
move in the same direction because the wage depends on both the production-labor
share and output per production worker.

\begin{proposition}[The wage and production-labor-share thresholds]
\label{prop:wage-share-thresholds}
Holding $K_t$ and $L_{Y,t}$ fixed, an expansion of AI services satisfies
\begin{align}
    \frac{\partial\log Y_t}{\partial\log X_t}
    &=(1-\alpha)s^X_t>0,
    \label{eq:output-ai-elasticity}\\
    \frac{\partial\log \ell^Y_t}{\partial\log X_t}
    &=-\left(1-\frac{1}{\sigma}\right)s^X_t,
    \label{eq:labor-share-ai-elasticity}\\
    \frac{\partial\log w_t}{\partial\log X_t}
    &=\left(\frac{1}{\sigma}-\alpha\right)s^X_t.
    \label{eq:wage-ai-elasticity}
\end{align}
Consequently, the production-labor share falls if and only if $\sigma>1$, whereas
the wage rises if and only if $\sigma<1/\alpha$. Since $\alpha\in(0,1)$, there is
a nonempty interval
\begin{equation}
    1<\sigma<\frac{1}{\alpha}
    \label{eq:wage-share-decoupling-region}
\end{equation}
in which an expansion of AI raises the wage while reducing production labor's share
of output.
\end{proposition}

\begin{proof}
The CES expenditure share in \eqref{eq:ai-composite-share} implies
\[
    \frac{\partial\log(1-s^X_t)}{\partial\log X_t}
    =
    -\left(1-\frac{1}{\sigma}\right)s^X_t.
\]
The output elasticity follows directly from
\eqref{eq:final-production}--\eqref{eq:service-composite}. Combining these two
expressions with \eqref{eq:labor-demand} yields
\eqref{eq:labor-share-ai-elasticity} and \eqref{eq:wage-ai-elasticity}.
\end{proof}

Because researchers and production workers receive the same wage, the aggregate
labor-income share is $w_tN_t/Y_t=(N_t/L_{Y,t})\ell^Y_t$. Conditional on a fixed
allocation of labor between production and research, its elasticity with respect to
$X_t$ is also given by \eqref{eq:labor-share-ai-elasticity}. Along an equilibrium
transition, however, changes in $H_t/N_t$ provide an additional margin. The
production-block threshold should therefore not be read as a complete
general-equilibrium result for the aggregate labor share.

Equation~\eqref{eq:wage-ai-elasticity} decomposes the wage response into a
productivity effect and a displacement effect:
\begin{equation}
    \frac{\partial\log w_t}{\partial\log X_t}
    =
    \underbrace{(1-\alpha)s^X_t}_{\text{productivity effect}}
    -
    \underbrace{\left(1-\frac{1}{\sigma}\right)s^X_t}
    _{\text{displacement effect}}.
    \label{eq:wage-decomposition}
\end{equation}
The first term is the increase in output per production worker. The second is the
decline in labor's income share. This decomposition maps the nested production
function into the productivity and displacement channels emphasized by
\citet{acemoglurestrepo2019}.

The two thresholds are not unique to this paper. \citet{yango2026} derives the same
conditions in a nested technology with physical and AI capital, while
\citet{autorkausik2026} establish the broader possibility of falling labor shares
and rising wages under constant returns to scale. Proposition
\ref{prop:wage-share-thresholds} serves here as a benchmark for the model's distinct
dynamic and market-structure mechanisms.

Allowing all quantities to change, the instantaneous wage growth rate obeys the
identity
\begin{equation}
    g_{w,t}
    =
    \alpha\left(g_{K,t}-g_{L_Y,t}\right)
    +
    s^X_t\left(\frac{1}{\sigma}-\alpha\right)
    \left(g_{X,t}-g_{L_Y,t}\right).
    \label{eq:dynamic-wage-decomposition}
\end{equation}
When the research-employment share is constant, $g_{L_Y,t}=n$. Capital deepening can
therefore raise wages even when the direct expansion of AI services lowers them.
Conversely, $\sigma<1/\alpha$ is not sufficient for wage growth along an equilibrium
transition if capital per production worker contracts. The monopoly problem below
determines the supply of $X_t$, while household saving and the resource constraint
determine $K_t$; the equilibrium wage response depends on both paths.

\subsubsection{Monopoly supply of AI services}

Holding $K_t$ and $L_{Y,t}$ fixed, the absolute elasticity of inverse demand faced by
the monopolist is
\begin{equation}
    \varepsilon^X_t
    \equiv
    -\frac{\partial\log p^X_t}{\partial\log X_t}
    =
    \frac{1-s^X_t}{\sigma}+\alpha s^X_t.
    \label{eq:inverse-demand-elasticity}
\end{equation}
For an interior solution with $\varepsilon^X_t<1$, the first-order condition is
\begin{equation}
    p^X_t(1-\varepsilon^X_t)=\frac{\xi}{A_t}.
    \label{eq:monopoly-ai-price}
\end{equation}
The markup and operating rent are
\begin{align}
    \mu^X_t
    &\equiv
    \frac{p^X_t}{\xi/A_t}
    =
    \frac{1}{1-\varepsilon^X_t}>1,
    \label{eq:monopoly-markup}\\
    \pi^X_t
    &=
    \varepsilon^X_t p^X_tX_t.
    \label{eq:monopoly-rents}
\end{align}
Conditional on $A_t$, $K_t$, and $L_{Y,t}$, the monopolist supplies less AI than an
allocation that prices services at marginal cost. The household receives the rents,
so the toy model retains the output distortion and the innovation incentive from
market power while abstracting from concentrated ownership.

\subsubsection{Research input choice and gradual automation}

The marginal products of the two research inputs in
\eqref{eq:effective-research} are
\begin{equation}
    E_{H,t}
    =(1-\nu)\left(\frac{E_t}{H_t}\right)^{1/\theta},
    \qquad
    E_{M,t}
    =\nu\left(\frac{E_t}{M_t}\right)^{1/\theta}.
    \label{eq:research-ces-marginal-products}
\end{equation}
The developer's first-order conditions are therefore
\begin{align}
    w_t
    &=
    q_t\chi\eta A_t^\phi E_t^{\eta-1}
    (1-\nu)\left(\frac{E_t}{H_t}\right)^{1/\theta},
    \label{eq:foc-human}\\
    \xi
    &=
    q_t\chi\eta A_t^\phi E_t^{\eta-1}
    \nu\left(\frac{E_t}{M_t}\right)^{1/\theta}.
    \label{eq:private-compute-foc}
\end{align}
Their ratio determines the research-input mix:
\begin{equation}
    \frac{M_t}{H_t}
    =
    \left[
        \frac{\nu}{1-\nu}\frac{w_t}{\xi}
    \right]^\theta.
    \label{eq:research-input-ratio}
\end{equation}

The unit cost of effective research services is
\begin{equation}
    P^E_t
    =
    \left[
        (1-\nu)^\theta w_t^{1-\theta}
        +\nu^\theta\xi^{1-\theta}
    \right]^{1/(1-\theta)}.
    \label{eq:research-unit-cost}
\end{equation}
Conditional input demands are
\begin{equation}
    H_t=(1-\nu)^\theta
    \left(\frac{P^E_t}{w_t}\right)^\theta E_t,
    \qquad
    M_t=\nu^\theta
    \left(\frac{P^E_t}{\xi}\right)^\theta E_t.
    \label{eq:research-conditional-demands}
\end{equation}
The developer can consequently choose the scale of research in two stages. For
$\eta\in(0,1)$, its optimal scale satisfies
\begin{equation}
    P^E_t=q_t\chi\eta A_t^\phi E_t^{\eta-1},
    \qquad
    E_t=
    \left(
        \frac{q_t\chi\eta A_t^\phi}{P^E_t}
    \right)^{1/(1-\eta)}.
    \label{eq:research-scale}
\end{equation}

Define the contribution of automated services to effective research as
\begin{equation}
    s^M_t
    \equiv
    \frac{\nu M_t^{(\theta-1)/\theta}}
    {(1-\nu)H_t^{(\theta-1)/\theta}
        +\nu M_t^{(\theta-1)/\theta}}.
    \label{eq:automated-research-share}
\end{equation}
At an interior cost-minimizing allocation, this is also the share of research
expenditure devoted to automated services:
\begin{equation}
    s^M_t=\frac{\xi M_t}{w_tH_t+\xi M_t}.
    \label{eq:automated-research-cost-share}
\end{equation}
Equations \eqref{eq:research-input-ratio} and
\eqref{eq:automated-research-share} imply
\begin{equation}
    \frac{s^M_t}{1-s^M_t}
    =
    \left(\frac{\nu}{1-\nu}\right)^\theta
    \left(\frac{w_t}{\xi}\right)^{\theta-1}.
    \label{eq:automation-odds}
\end{equation}
More generally, if the unit cost of automated research varies at rate $g_{\xi,t}$,
the automation share evolves according to
\begin{equation}
    \dot s^M_t
    =
    s^M_t(1-s^M_t)(\theta-1)
    \left(g_{w,t}-g_{\xi,t}\right).
    \label{eq:automation-share-dynamics}
\end{equation}
The baseline holds $\xi$ constant, so $g_{\xi,t}=0$.

\begin{proposition}[Gradual research automation]
\label{prop:gradual-automation}
Suppose $\theta>1$. If $w_t/\xi\to\infty$, then $s^M_t\to1$ and
$H_t/E_t\to0$. If $w_t/\xi$ remains bounded, $\theta>1$ alone does not imply that
the human contribution disappears.
\end{proposition}

Proposition~\ref{prop:gradual-automation} replaces a discrete automation threshold
with a continuous equilibrium transition. The research elasticity determines
whether human research can become asymptotically inessential. The production
elasticity $\sigma$ affects whether that limit is reached by determining the wage
path. The private shadow value continues to satisfy
\begin{equation}
    (r_t-\delta_K)q_t
    =
    \pi^X_{A,t}+q_tF_{A,t}+\dot q_t,
    \qquad
    F_{A,t}=\phi\frac{\dot A_t}{A_t}.
    \label{eq:private-costate}
\end{equation}

\subsubsection{Human-dominated research}

When $s^M_t$ is close to zero,
\begin{equation}
    E_t
    \simeq
    (1-\nu)^{\theta/(\theta-1)}H_t.
    \label{eq:human-dominated-research}
\end{equation}
If a constant share $s_H\in(0,1)$ of the population works in research,
\begin{equation}
    H_t=s_HN_t,
    \qquad
    L_{Y,t}=(1-s_H)N_t,
    \label{eq:research-share}
\end{equation}
the ideas equation implies
\begin{equation}
    g_A^H
    =
    \frac{\eta n}{1-\phi}.
    \label{eq:human-dominated-ga}
\end{equation}
This is a semi-endogenous benchmark for the early, human-dominated phase. It is not
generally a permanent balanced-growth path once automated research is feasible. If
the wage grows relative to $\xi$, equation \eqref{eq:automation-share-dynamics}
moves the economy away from this phase.

\subsubsection{AI-dominated research dynamics}

When $s^M_t\to1$, the research aggregator satisfies
\begin{equation}
    E_t
    \simeq
    \nu^{\theta/(\theta-1)}M_t.
    \label{eq:ai-dominated-research}
\end{equation}
The asymptotic law of motion is therefore
\begin{equation}
    \dot A_t
    =
    \kappa A_t^\phi M_t^\eta,
    \qquad
    \kappa\equiv
    \chi\nu^{\eta\theta/(\theta-1)}.
    \label{eq:ai-dominated-ideas}
\end{equation}
Its growth rate obeys
\begin{equation}
    \frac{\dot g_{A,t}}{g_{A,t}}
    =
    (\phi-1)g_{A,t}+\eta g_{M,t}.
    \label{eq:ai-growth-rate-dynamics}
\end{equation}
If $g_{M,t}$ converges to a positive constant and $\phi<1$, a constant capability
growth rate must satisfy
\begin{equation}
    g_A^\infty
    =
    \frac{\eta g_M^\infty}{1-\phi}.
    \label{eq:ai-dominated-ga}
\end{equation}
Research growth is no longer tied directly to population growth. It is tied to the
growth of automated-research services, whose equilibrium path depends on output,
profits, and the private value of capability.

\subsubsection{A closed-form AI-dominated path under Cobb--Douglas production}

The Cobb--Douglas limit of the production CES provides a closed-form benchmark.
Define
\begin{equation}
    \beta\equiv(1-\alpha)\omega,
    \qquad
    \upsilon\equiv\frac{\beta}{1-\alpha-\beta}.
    \label{eq:beta-definition}
\end{equation}
Production becomes
\begin{equation}
    Y_t=K_t^\alpha L_{Y,t}^{1-\alpha-\beta}X_t^\beta.
    \label{eq:cd-production}
\end{equation}
The monopolist chooses
\begin{align}
    p^X_t&=\frac{\xi}{\beta A_t},
    &
    U_t&=\frac{\beta^2}{\xi}Y_t,
    &
    X_t&=\frac{\beta^2A_t}{\xi}Y_t,
    \label{eq:cd-monopoly}\\
    \pi^X_t&=\beta(1-\beta)Y_t.
    \label{eq:cd-monopoly-rents}
\end{align}
Substitution into production gives
\begin{equation}
    Y_t
    =
    \left(\frac{\beta^2}{\xi}\right)^{\beta/(1-\beta)}
    A_t^{\beta/(1-\beta)}
    K_t^{\alpha/(1-\beta)}
    L_{Y,t}^{(1-\alpha-\beta)/(1-\beta)}.
    \label{eq:reduced-production}
\end{equation}

Let $\gamma^\infty$ denote the common growth rate of output, capital,
consumption, and inference compute per capita. Along an AI-dominated path with a
constant research-resource share $m_R\equiv\xi M_t/Y_t$,
\begin{equation}
    \gamma^\infty=\upsilon g_A^\infty,
    \qquad
    g_M^\infty=n+\gamma^\infty.
    \label{eq:cd-ai-growth-relations}
\end{equation}
Combining \eqref{eq:ai-dominated-ga} and
\eqref{eq:cd-ai-growth-relations} yields
\begin{align}
    g_A^\infty
    &=
    \frac{\eta n}
    {1-\phi-\eta\upsilon},
    \label{eq:cd-ai-ga}\\
    \gamma^\infty
    &=
    \frac{\upsilon\eta n}
    {1-\phi-\eta\upsilon}.
    \label{eq:cd-ai-gamma}
\end{align}
A positive constant-growth path requires
\begin{equation}
    1-\phi-\eta\upsilon>0.
    \label{eq:cd-ai-bg-condition}
\end{equation}
If this condition fails, the feedback from output to automated research and from
research back to output is too strong for a positive constant-growth solution. The
failure of \eqref{eq:cd-ai-bg-condition} identifies a candidate acceleration regime;
the transition system determines whether research and output accelerate or instead
move to a boundary allocation.

The remaining balanced-growth objects are
\begin{align}
    g_Y=g_K=g_C=g_U=g_M&=n+\gamma^\infty,
    \label{eq:aggregate-bg-rates}\\
    g_X&=n+\gamma^\infty+g_A^\infty,
    &
    g_w&=\gamma^\infty,
    &
    g_{p^X}&=-g_A^\infty,
    \label{eq:other-bg-rates}\\
    r^\infty&=\rho+\delta_K+\gamma^\infty,
    &
    \frac{K_t}{Y_t}
    &=
    \frac{\alpha}{\rho+\delta_K+\gamma^\infty}.
    \label{eq:bg-capital-output}
\end{align}
The costate equation and the research first-order condition imply
\begin{equation}
    m_R^\infty
    =
    \frac{\beta^2\eta g_A^\infty}
    {\rho-n+(1-\phi)g_A^\infty}.
    \label{eq:bg-research-resource-share}
\end{equation}
Consequently,
\begin{align}
    \frac{I_t}{Y_t}
    &=
    \frac{\alpha(n+\gamma^\infty+\delta_K)}
    {\rho+\delta_K+\gamma^\infty},
    \label{eq:bg-investment-share}\\
    \frac{C_t}{Y_t}
    &=
    1-
    \frac{\alpha(n+\gamma^\infty+\delta_K)}
    {\rho+\delta_K+\gamma^\infty}
    -\beta^2-m_R^\infty.
    \label{eq:bg-consumption-share}
\end{align}
An interior path requires a positive consumption share.

Because $g_w=\gamma^\infty>0$ and $\xi$ is constant,
Proposition~\ref{prop:gradual-automation} applies. More precisely,
\begin{equation}
    g_{H/N}
    =
    -(\theta-1)\gamma^\infty<0.
    \label{eq:human-research-population-share}
\end{equation}
The fraction of the population employed in frontier research therefore converges to
zero. The number of human researchers need not fall in levels: it grows at rate
$n-(\theta-1)\gamma^\infty$ and declines only if the second term exceeds population
growth.

\subsubsection{Why the production CES matters}

The Cobb--Douglas case is the knife-edge at which expenditure shares remain
constant. For $\sigma\neq1$, the monopoly condition
\eqref{eq:monopoly-ai-price} implies different asymptotic regimes as the marginal
cost $\xi/A_t$ falls.

\begin{proposition}[Production regimes under increasingly capable AI]
\label{prop:production-ces-regimes}
Suppose $A_t\to\infty$ and the service-market solution remains interior.
\begin{enumerate}[label=(\roman*)]
    \item If $\sigma<1$, the monopolist's AI share converges to
    \begin{equation}
        \bar s^X
        =
        \frac{1-\sigma}{1-\alpha\sigma}.
        \label{eq:low-sigma-ai-share}
    \end{equation}
    The ratio $X_t/L_{Y,t}$ remains bounded and the productive effect of further
    reductions in $\xi/A_t$ vanishes.
    \item If $\sigma=1$, factor shares are constant and the AI-dominated
    balanced-growth path is characterized by
    \eqref{eq:cd-ai-ga}--\eqref{eq:bg-consumption-share}.
    \item If $\sigma>1$, $s^X_t\to1$, the production-labor share converges to zero,
    and
    \begin{equation}
        \frac{Y_t}{K_t}
        \sim
        D_\sigma A_t^{(1-\alpha)/\alpha},
        \qquad D_\sigma>0.
        \label{eq:high-sigma-ak-limit}
    \end{equation}
    A conventional balanced-growth path with $g_A>0$ cannot have a constant
    capital--output ratio. The return to capital and the aggregate growth rate rise
    with capability.
\end{enumerate}
\end{proposition}

For $\sigma<1$, a rising wage relative to $\xi$ is not guaranteed: labor remains
essential and the monopoly supply of AI services approaches a finite
revenue-maximizing ratio. Research automation can therefore stop short of
$s^M=1$ even though $\theta>1$. For $\sigma=1$, the positive wage growth in
\eqref{eq:other-bg-rates} makes automated research asymptotically dominant. For
$\sigma>1$, production itself becomes an acceleration mechanism. Research
automation reinforces that mechanism, but the wage path must be solved jointly with
capital accumulation before concluding that the human research share disappears.

\subsubsection{Private and social research}

The planner's first-order condition for research compute is
\begin{equation}
    Q_tF_{M,t}=\xi.
    \label{eq:social-compute-foc}
\end{equation}
The private condition is $q_tF_{M,t}=\xi$. At the private allocation, the
uninternalized marginal social benefit is therefore
\begin{equation}
    Q_tF_{M,t}-\xi=(Q_t-q_t)F_{M,t}.
    \label{eq:knowledge-wedge}
\end{equation}

\begin{proposition}[The knowledge and appropriability wedge]
\label{prop:knowledge-wedge}
If $Q_t>q_t$, the integrated monopolist invests too little in research compute at the
margin relative to the social allocation.
\end{proposition}

There are thus two distinct distortions. Static market power restricts current AI
services, while incomplete appropriation of consumer surplus and knowledge benefits
can depress frontier research. Price regulation alone may correct the first
distortion while weakening the rents that finance research. This is the central
policy trade-off developed below.

\input{sections/04_extended_model}
\input{sections/05_regulation}
\section{Conclusion}
\label{sec:conclusion}

The economic consequences of advanced AI depend on more than whether machines can
perform existing jobs. A distinct transition occurs when AI can perform the research
required to produce better AI. The toy model represents this transition as a gradual
change in the composition of research. When human and automated-research services are
gross substitutes, the automated share rises with the research wage relative to its
unit resource cost.

Population growth matters while human research dominates. With $\phi<1$ and a
constant research-employment share, capability growth is semi-endogenous and equals
$\eta n/(1-\phi)$. When automated research dominates, capability growth is tied
instead to the growth of automated-research services. In the Cobb--Douglas production
benchmark, a constant-growth path exists if
$1-\phi-\eta\beta/(1-\alpha-\beta)>0$. If this condition fails, the feedback between
output and automated research is too strong for a positive constant-growth solution.

The production block also separates the wage response from the response of the
production-labor share. An expansion of AI services reduces this share when
$\sigma>1$, but it raises the wage when $\sigma<1/\alpha$. Hence both outcomes occur
when $1<\sigma<1/\alpha$. This conditional result is a benchmark, not the paper's
standalone contribution. In general equilibrium, capital deepening and the
reallocation of workers between production and research also affect wages and the
aggregate labor share.

The production elasticity determines whether the economy reaches that asymptotic
research regime. When $\sigma<1$, labor remains essential and the productive return
to further reductions in the unit cost of AI services vanishes. When $\sigma=1$,
positive wage growth makes automated research asymptotically dominant. When
$\sigma>1$, the production-labor share converges to zero and the return to capital
rises with capability, so a conventional balanced-growth path with positive
capability growth cannot maintain a constant capital--output ratio.

Market power creates a genuine policy trade-off. The monopolist restricts the current
supply of AI services, but its operating rents raise the private return to frontier
research. Because it does not appropriate all consumer surplus and knowledge
benefits, private research remains insufficient whenever $Q_t>q_t$. Regulating the
AI-service price without replacing the lost innovation return can therefore improve
current deployment while slowing capability growth. The natural policy comparison is
between laissez-faire monopoly, regulated access, explicit research support, and a
joint policy that targets both margins.


\bibliographystyle{apalike}
\bibliography{references}

\appendix
\section{Proofs}

\begin{proof}[Proof of Proposition~\ref{prop:gradual-automation}]
Equation~\eqref{eq:automation-odds} gives
\[
    \frac{s^M_t}{1-s^M_t}
    =
    \left(\frac{\nu}{1-\nu}\right)^{\sigma_{HM}}
    \left(\frac{w_t}{\xi}\right)^{\sigma_{HM}-1}.
\]
If $\sigma_{HM}>1$ and $w_t/\xi\to\infty$, the right-hand side diverges and
$s^M_t\to1$. Moreover, the unit cost in
\eqref{eq:research-unit-cost} converges to
\[
    P^E_t\longrightarrow
    \nu^{-\sigma_{HM}/(\sigma_{HM}-1)}\xi.
\]
The first expression in \eqref{eq:research-conditional-demands} then implies
$H_t/E_t\to0$. If $w_t/\xi$ remains bounded, the odds in
\eqref{eq:automation-odds} remain bounded, so $\sigma_{HM}>1$ does not by itself imply
$s^M_t\to1$.
\end{proof}

\begin{proof}[Proof of Proposition~\ref{prop:singularity-conditions}]
Let $z_t\equiv1/g_{A,t}$. Equation
\eqref{eq:ces-transition-growth-dynamics} implies the linear equation
\[
    \dot z_t=-\eta n z_t+D(s^M_t).
\]
Multiplying by $e^{\eta nt}$ and integrating gives
\eqref{eq:time-varying-singularity-condition}. If $D(s^M_t)\geq0$, the bracket
in that equation remains positive. If $D>0$ is constant, the logistic equation for
$g_A$ has the stable positive fixed point $\eta n/D$.

If $D=0$, direct integration gives
\[
    g_{A,t}=g_{A,0}e^{\eta nt},
    \qquad
    A_t
    =
    A_0\exp\left[
        \frac{g_{A,0}}{\eta n}
        \left(e^{\eta nt}-1\right)
    \right].
\]
Both expressions are finite for every finite $t$. If $D<0$ is constant, the
solution for the growth rate is
\[
    g_{A,t}
    =
    \frac{g_{A,0}e^{\eta nt}}
    {1-\frac{(-D)g_{A,0}}{\eta n}
        \left(e^{\eta nt}-1\right)}.
\]
The denominator reaches zero at the date in
\eqref{eq:constant-feedback-singularity-time}. Near that date, $g_A$ is
proportional to $(T^*-t)^{-1}$, so both $g_A$ and $A$ diverge. For a time-varying
$D(s^M_t)$, the same divergence occurs exactly when the bracket in
\eqref{eq:time-varying-singularity-condition} reaches zero.
\end{proof}

\begin{proof}[Proof of Proposition~\ref{prop:production-ces-regimes}]
Let $x_t\equiv X_t/L_{Y,t}$, $y_t\equiv Y_t/L_{Y,t}$, and write
$s_t\equiv s^X_t$. The monopoly condition can be written as
\[
    \frac{\xi x_t}{A_ty_t}
    =
    (1-\alpha)s_t[1-\varepsilon(s_t)],
    \qquad
    \varepsilon(s)
    =
    \frac{1-s}{\sigma_{XL}}+\alpha s.
\]

For $\sigma_{XL}<1$, revenue approaches zero both as $x_t\to0$ and as
$x_t\to\infty$. As marginal cost converges to zero, the monopolist approaches the
interior revenue-maximizing quantity, which satisfies $\varepsilon(\bar s^X)=1$.
Solving this condition gives
\[
    \bar s^X=\frac{1-\sigma_{XL}}{1-\alpha\sigma_{XL}}.
\]
Because the CES mapping between $s^X$ and $X/L_Y$ is one-to-one, the corresponding
service ratio is finite.

For $\sigma_{XL}=1$, the CES contribution $s^X_t$ equals the constant weight $\omega$.

For $\sigma_{XL}>1$, revenue is unbounded as $x_t\to\infty$ when marginal cost is zero,
so $x_t\to\infty$ and $s^X_t\to1$. Consequently,
$\varepsilon^X_t\to\alpha$, and the monopoly first-order condition implies
\[
    \frac{\xi U_t}{Y_t}\longrightarrow(1-\alpha)^2.
\]
Since
\[
    Z_t\sim
    \omega^{\sigma_{XL}/(\sigma_{XL}-1)}X_t,
\]
substituting $X_t=A_tU_t$ and the limiting compute share into production yields
\[
    \frac{Y_t}{K_t}
    \sim
    \left[
        \frac{\omega^{\sigma_{XL}/(\sigma_{XL}-1)}(1-\alpha)^2}{\xi}
    \right]^{(1-\alpha)/\alpha}
    A_t^{(1-\alpha)/\alpha}.
\]
This establishes \eqref{eq:high-sigma-ak-limit}.
\end{proof}

\begin{proof}[Derivation of Equation~\eqref{eq:bg-research-resource-share}]

In the Cobb--Douglas production benchmark, the envelope theorem and
\eqref{eq:cd-monopoly} imply
\[
    \pi^X_{A,t}=\frac{\beta^2Y_t}{A_t}.
\]
The research first-order condition in the AI-dominated limit gives
\[
    \xi=q_t\eta\frac{\dot A_t}{M_t},
    \qquad
    \frac{q_tA_t}{Y_t}
    =
    \frac{m_R^\infty}{\eta g_A^\infty}.
\]
Because the right-hand side is constant on a balanced-growth path,
$g_q=g_Y-g_A$. Dividing \eqref{eq:private-costate} by $q_t$, using
$r-\delta=\rho+\gamma^\infty$ and $g_Y=n+\gamma^\infty$, and substituting the
two expressions above gives
\[
    \rho-n
    =
    \frac{\beta^2\eta g_A^\infty}{m_R^\infty}
    -(1-\phi)g_A^\infty.
\]
Solving for $m_R^\infty$ yields
\eqref{eq:bg-research-resource-share}.
\end{proof}

\begin{proof}[Proof of Proposition~\ref{prop:knowledge-wedge}]
Evaluate the planner's marginal net benefit from research compute at the private
allocation. Using \eqref{eq:private-compute-foc},
\[
    Q_tF_{M,t}-\xi
    =
    (Q_t-q_t)F_{M,t}.
\]
If $Q_t>q_t$, the expression is positive and the planner increases research compute
locally relative to the private allocation.
\end{proof}


\end{document}
